\chapter{Related Work}\label{chap:relatedWork}

An early use of ILPs in planning was by~\cite{Bylander1997LplanHeuristic}, who developed the Lplan heuristic, the first admissible heuristic for partial-order planning. This heuristic was developed from an ILP method which can solve any classical planning problem. This technique utilises the fact that classical planning is NP-complete for plans of polynomially bounded length, as shown in Corollary 1 of the paper by~\cite{Turner2002PolynomialBoundedPlans}. To find the optimal plan length, an ILP is created which solves the planning problem with a bound of 1, and if it is infeasible a new ILP is created for a bound of 2, with this process repeated until the first feasible bound is found, which hence gives the minimum plan length. However, this formulation is inefficient as a direct method for solving planning problems, and is additionally problematic as it will run forever on infeasible planning problems, as it has no method for detecting infeasibility. Instead, the ILP is relaxed to an LP for use as a heuristic. As an LP allows non-integer variable values, this allows for strange solutions, such as half of two different actions being executed at one time step. It also still contains efficiency concerns due to having to run the LP multiple times for each possible bound on solution length. However, the main advantage of the Lplan heuristic is that it requires many fewer node visits than other heuristics of the time.\\


While future papers on the use of ILPs for planning shifted away from partial-order planning as it became less popular, there were still many papers focused on modeling classical planning problems, and relaxations of those problems, as ILPs, as well as the use of these ILPs as heuristics. Many early attempts at using ILPs for planning aimed to solve planning problems directly with ILPs, using similar techniques to~\cite{Bylander1997LplanHeuristic}, where successive ILPs were used to check each possible bound until a solution was found. The culmination of this type of ILP solver was the Optiplan model by~\cite{VanDenBriel2005Optiplan}. Optiplan combined the \textit{state-change formulation} method of~\cite{Vossen1999ILPModelsInPlanning}, which used variables to directly express the addition, deletion, and persistence of propositions over time, with additional refinement strategies developed by~\cite{Dimopoulos2014ImprovedILPModelsPlanning}. However, SAT planning methods, which used a similar technique of iterative checking of higher bounds on plan size, but using satisfiability formulas instead of ILPs, continued to outperform ILP based methods, and so research into ILP based planners slowed.\\


Despite the lack of promise in the field of ILP planners, the use of LPs for creating admissible heuristics, as with the Lplan heuristic, still remained potentially promising. An admissible heuristic using an ILP model was developed by~\cite{VanDenBriel2007LpHeuristicPlanning}, which completely relaxes all ordering constraints, including allowing the effects of one actions to be used as preconditions for an earlier action, to create an ILP formulation. This formulation is then relaxed to an LP and used as a heuristic. This heuristic was shown to outperform Lplan on many planning instances. The SEQ Heuristic, introduced by~\cite{Bonet2013AdmissibleLPHeuristic} uses directed bipartitite graph relaxations of planning problems to generate relaxed solutions which provide a minimal number of each action that must be executed to real a goal state. By counting the total of all of these actions, an admissible heuristic is found. This calculation requires an ILP, but using its LP relaxation as a heuristic was found to be competitive with state of the art optimal planning heuristics. A third LP based heuristic was a method for representing pattern database heuristics developed by~\cite{Pommerening2013PatternDatabaseLPHeuristics}. Pattern database heuristics store a large number of perfect cost estimates for small subproblems, and combine these subproblems together to estimate the cost of larger problems. An important component of pattern database heuristics is how to determine which patterns are relevant. This paper presents the Post-Hoc Optimisation heuristic, which links patterns that modify certain variables to actions which modify the same variables, joining them into operator groups and obtaining bounds on the minimum number of actions from each group required for a solution. This heuristic outperformed previous pattern database heuristic methods. You may note that all three of the heuristics discussed utilise some form of counting the number of each action required for some lower bound on solution length. This was also recognised by~\cite{Pommerening2014LPHeuristics}, who developed a more general \textit{operator counting} heuristic, which was a framework for developing LPs which could represent any of the prior three heuristics, as well as other similar heuristics. It could also combine multiple heuristics that fit into the operator counting heuristic together, allowing it to dominate all individual operator counting heuristics by running them all at once and taking the maximum heuristic value. This technique was found to be very powerful, and although it unfortunately cannot be used to entirely encode the heuristics we present in this paper, it is utilised to optimise our ILP model (TODO: in future check this again, as I may not end up having time to implement operator counting optimisations).\\


Another use case, using ILPs to compute the standard delete-relaxation heuristic of a state-based planning problem ($h^+$), was developed by~\cite{Imai2015ILPDeleteRelaxedPlanning}. The LP relaxation was shown to always produce integer results, and was competitive with other methods of computing the $h^+$ heuristic. This heuristic was used as the basis of a heiarchical task network (HTN) planning heuristic by~\cite{Hoeller2020ILP}, which computed a delete and ordering relaxed (DOR) HTN problem with an ILP as a heuristic. What is particularly relevant about this paper is that it also features a type of delete relaxation which, while polytime in state-based planning, is NP-complete in HTN planning, as is the case for POCL planning. The DOR heuristic was more informed than any other HTN heuristics in the literature, and remained competitive despite being NP-complete. The LP relaxation of the heuristic also performed well.\\


Within the field of POCL planning, ~\cite{Bercher2013POCLHeuristicsViaEncoding} developed the only heuristic which accounts for all POCL search node information. This was done by translating a partial plan into a classical planning problem, which could then be solved as a state-based plan, using state-based relaxations and heuristics. This allowed for state-based admissible heuristics to be used in POCL planning, where the only prior admissible heuristic for partial order planning was the Lplan heuristic. This heuristic was slow compared to heuristics built specifically for POCL, with much of this being due to slow translation times, and unoptimised translations of causal links, an issue we solve in our heuristic. However, it did show promise for potential reduction in the number of explored search nodes. The theoretical background for our heuristic comes from the research by~\cite{Bercher2021POCLComplexities}, which in Table 1 lists the complexity of many different relaxation possibilities for POCL planning. It finds that by only relaxing the delete effects of insertable actions; and not relaxing any actions, causal links, and ordering constraints within a partial plan; an NP-complete heuristic can be developed. Hence, such a heuristic is computable with an ILP model. The paper also explores the potential advantages of such a method over polytime POCL heuristics.
